% CHAPTER 2 - Add this after Chapter 1 in your main .tex file
% Or % CHAPTER 2 - Add this after Chapter 1 in your main .tex file
% Or % CHAPTER 2 - Add this after Chapter 1 in your main .tex file
% Or % CHAPTER 2 - Add this after Chapter 1 in your main .tex file
% Or \input{chapter2.tex} from your main document

\chapter{SYSTEM COMPONENTS AND WORKING PRINCIPLES}

\section{Overview}

The design and implementation of residential photovoltaic (PV) systems---whether on-grid or hybrid---require an in-depth understanding of the individual components, their operating principles, and how they interact. This chapter presents a detailed examination of the core components: solar PV modules, inverters, battery energy storage systems, and grid connection mechanisms.


\section{Solar Photovoltaic (PV) Modules}

\subsection{Operating Principle: The Photovoltaic Effect}

Solar PV modules convert incident solar radiation directly into electrical energy through the photovoltaic effect---a quantum mechanical process first observed by Becquerel in 1839 and first practically demonstrated by Bell Laboratories in 1954 \cite{chapin1954}.

The photovoltaic effect occurs within semiconductor materials---most commonly crystalline silicon (Si)---engineered to contain a p-n junction. When photons strike the semiconductor surface with energy greater than the material's bandgap energy (approximately 1.12~eV for silicon), they excite valence electrons into the conduction band, creating electron-hole pairs. The internal electric field at the p-n junction separates these charge carriers: electrons are swept toward the n-type (phosphorus-doped) layer, while holes migrate toward the p-type (boron-doped) layer. This charge separation establishes a potential difference across the cell, and when an external circuit is connected, current flows \cite{green2003,luque2011}.

\begin{figure}[H]
    \centering
    \includegraphics[width=0.75\textwidth]{figures/fig2_1_pv_conversion.png}
    \caption{The conversion of solar radiation to electricity in PV systems, showing the p-n junction, photon absorption, electron flow, and hole flow mechanisms.}
    \label{fig:pv_conversion}
\end{figure}

The performance of a PV cell is characterized by several key electrical parameters measured under Standard Test Conditions (STC: 1000~W/m$^{2}$ irradiance, 25$^{\circ}$C cell temperature, and an Air Mass of 1.5):

\begin{itemize}
    \item \textbf{Open-Circuit Voltage ($V_{oc}$):} This represents the maximum voltage available from a solar cell when no current is flowing, typically ranging from 0.5 to 0.7~V for standard silicon cells.
    \item \textbf{Short-Circuit Current ($I_{sc}$):} This is the maximum current output at zero voltage, which is directly proportional to the incident solar irradiance and the active cell area.
    \item \textbf{Maximum Power Point ($P_{max}$):} During normal operation, the cell operates at this specific voltage and current point where the electrical power output is maximized.
    \item \textbf{Fill Factor (FF):} Defined as the ratio of $P_{max}$ to the product of $V_{oc}$ and $I_{sc}$. High-quality cells typically exhibit a fill factor between 0.75 and 0.85.
    \item \textbf{Conversion Efficiency ($\eta$):} This dictates the proportion of incident solar power that is successfully converted into usable electrical power.
\end{itemize}

The fundamental relationships are:

\begin{equation}
    P_{max} = V_{mp} \times I_{mp} = FF \times V_{oc} \times I_{sc}
    \label{eq:pmax}
\end{equation}

\begin{equation}
    \eta = \frac{P_{max}}{G \times A} \times 100\%
    \label{eq:efficiency}
\end{equation}

\noindent where $G$ is solar irradiance (W/m$^{2}$) and $A$ is cell area (m$^{2}$).

\begin{figure}[H]
    \centering
    \includegraphics[width=0.7\textwidth]{figures/fig2_2_iv_curve.png}
    \caption{Typical I-V and P-V characteristic curves of a solar cell showing open-circuit voltage ($V_{oc}$), short-circuit current ($I_{sc}$), and the maximum power point ($P_{max}$).}
    \label{fig:iv_curve}
\end{figure}


\subsection{PV Module Construction and Types}

A PV module consists of multiple cells electrically connected in series and parallel to achieve desired voltage and current ratings. Cells are encapsulated between a front glass cover and rear backsheet with EVA encapsulant, providing mechanical protection and moisture resistance.

\begin{table}[H]
\centering
\caption{Comparison of Major PV Module Technologies}
\label{tab:pv_comparison}
\small
\begin{tabular}{@{}p{3.8cm} c c c@{}}
\toprule
\textbf{Parameter} & \textbf{Monocrystalline Si} & \textbf{Polycrystalline Si} & \textbf{Thin-Film} \\
\midrule
Cell Efficiency & 18--24\% & 16--20\% & 10--13\% \\
Module Efficiency & 17--22\% & 15--18\% & 9--12\% \\
Temp.\ Coefficient & --0.3 to --0.4 \%/$^{\circ}$C & --0.4 to --0.5 \%/$^{\circ}$C & --0.2 to --0.3 \%/$^{\circ}$C \\
Degradation Rate & 0.3--0.5 \%/yr & 0.5--0.7 \%/yr & 0.5--1.0 \%/yr \\
Lifespan & 25--30+ years & 25--30 years & 20--25 years \\
Space Efficiency & High & Medium & Low \\
Cost per Watt & Moderate--High & Low--Moderate & Low \\
\bottomrule
\end{tabular}
\end{table}

Monocrystalline silicon modules dominate the residential market due to superior efficiency and favorable temperature coefficients. The transition from PERC to TOPCon and HJT architectures has further improved efficiencies while reducing degradation \cite{longi2024,green2022}.


\subsection{Environmental Factors Affecting Performance}

The actual field output of a PV module is highly dependent on localized environmental parameters that cause deviations from STC ratings:

\begin{itemize}
    \item \textbf{Operating Temperature:} Module efficiency exhibits a negative correlation with increasing cell temperature, typically degrading at a rate of --0.3 to --0.5~\%/$^{\circ}$C. In Erbil's extreme summer climate, where rooftop module temperatures can reach 60--70$^{\circ}$C, temperature-induced power losses can reach 15--20\% relative to baseline STC values \cite{skoplaki2009}.
    \item \textbf{Soiling and Dust Accumulation:} The accumulation of airborne particulate matter poses a significant challenge. Studies in arid Middle Eastern climates indicate soiling can induce yield losses of 15--30\% if regular maintenance is neglected, making periodic cleaning essential during the Kurdistan Region's dry summers \cite{mani2010}.
    \item \textbf{Wind Speed:} While moderate wind velocities can marginally improve performance by providing convective cooling to the modules, extreme wind loads necessitate robust structural mounting to prevent mechanical failure.
    \item \textbf{Shading:} Even partial shading from nearby structures or vegetation can disproportionately reduce total array output. Because cells within a module are series-connected, a shaded cell acts as a reverse-biased resistor, restricting current flow and creating localized hot spots.
\end{itemize}

\begin{figure}[H]
    \centering
    \includegraphics[width=0.75\textwidth]{figures/fig2_3_rooftop_orientation.png}
    \caption{Typical residential rooftop solar PV installation showing panel arrangement, mounting structure, tilt angle, and exposure to sunlight with minimal shading conditions.}
    \label{fig:rooftop}
\end{figure}


\section{Inverters}

\subsection{Function and Operating Principle}

The primary function of a solar inverter is the conversion of direct current (DC) electricity generated by the PV array into alternating current (AC). However, modern inverters act as the central intelligence hub of the PV system, executing several critical ancillary functions:

\begin{itemize}
    \item \textbf{Maximum Power Point Tracking (MPPT):} The inverter continuously samples the voltage and current of the PV array and adjusts the electrical operating point to ensure maximum power extraction under fluctuating conditions, frequently achieving tracking efficiencies exceeding 99\% \cite{esram2007}.
    \item \textbf{Grid Synchronization:} For grid-tied operations, the inverter precisely matches the output voltage, phase, and frequency to the local utility specifications (50~Hz in Iraq) while maintaining Total Harmonic Distortion (THD) below 3--5\%.
    \item \textbf{Anti-Islanding Protection:} Crucially for safety, grid-interactive inverters must rapidly detect utility outages and automatically disconnect from the grid per IEEE~1547 standards \cite{ieee1547}, preventing power back-feed into compromised distribution lines.
    \item \textbf{System Monitoring:} Contemporary inverters integrate robust data logging and communication protocols (Wi-Fi, Ethernet, RS485), enabling real-time system monitoring, fault diagnostics, and remote management.
\end{itemize}


\subsection{Types of Inverters}

\begin{table}[H]
\centering
\caption{Comparison of Inverter Types for Residential PV Systems}
\label{tab:inverter_comparison}
\small
\begin{tabular}{@{}p{3.5cm} c c c@{}}
\toprule
\textbf{Parameter} & \textbf{String} & \textbf{Microinverter} & \textbf{Hybrid} \\
\midrule
Peak Efficiency & 96--98\% & 95--97\% & 93--97\% \\
MPPT Level & System (1--2) & Module-level & System + Battery \\
Shading Tolerance & Low & High & Low--Medium \\
Battery Integration & No & No & Yes (built-in) \\
Islanding Capability & No & No & Yes \\
Cost per Watt & Lowest & Highest & Moderate--High \\
Typical Lifespan & 10--15 yr & 15--25 yr & 10--15 yr \\
\bottomrule
\end{tabular}
\end{table}

For this study, a hybrid inverter is employed for the hybrid system, while a standard string inverter is utilized for the on-grid system.


\section{Battery Energy Storage Systems (BESS)}

\subsection{Role of Battery Storage}

Within a hybrid PV architecture, Battery Energy Storage Systems (BESS) are utilized to resolve the fundamental temporal mismatch between peak solar generation and residential consumption, serving three primary strategic functions:

\begin{itemize}
    \item \textbf{Energy Time-Shifting:} By storing excess diurnal generation for nocturnal use, households can increase their volumetric self-consumption of solar energy from a baseline of 30--40\% to upwards of 60--80\% \cite{koskela2019}.
    \item \textbf{Backup Power:} In the context of the Kurdistan Region's unreliable electricity infrastructure, batteries provide uninterruptible power supply during sudden utility grid outages.
    \item \textbf{Peak Shaving:} Intelligent hybrid inverters can utilize stored battery capacity to deliberately discharge during periods of high grid demand, minimizing reliance on expensive utility tariffs or neighborhood diesel generator subscriptions.
\end{itemize}

\begin{figure}[H]
    \centering
    \includegraphics[width=0.8\textwidth]{figures/fig2_4_wiring_diagram.png}
    \caption{Electrical wiring and connections in a residential PV system showing integration between solar panels, charge controller, battery, inverter, and load.}
    \label{fig:wiring}
\end{figure}


\subsection{Battery Technologies}

\begin{table}[H]
\centering
\caption{Comparison of Battery Technologies for Residential PV Systems}
\label{tab:battery_comparison}
\small
\begin{tabular}{@{}p{3.8cm} c c c@{}}
\toprule
\textbf{Parameter} & \textbf{Lead-Acid} & \textbf{Li-Ion (NMC)} & \textbf{LiFePO4 (LFP)} \\
\midrule
Energy Density & 30--50 Wh/kg & 150--220 Wh/kg & 90--160 Wh/kg \\
Cycle Life (80\% DoD) & 300--800 & 2,000--4,000 & 3,000--10,000 \\
Recommended DoD & 50\% & 80--90\% & 80--100\% \\
Round-Trip Efficiency & 75--85\% & 90--95\% & 92--98\% \\
Calendar Life & 3--5 yr & 8--12 yr & 10--15 yr \\
Safety & H$_2$ off-gassing & Thermal runaway risk & Excellent stability \\
Cost per kWh (2024) & \$100--200 & \$250--400 & \$200--350 \\
\bottomrule
\end{tabular}
\end{table}

LFP chemistry is increasingly preferred for residential applications due to excellent thermal stability, non-toxicity, and extended cycle life \cite{zakeri2015,hesse2017}.


\subsection{Battery Sizing}

The required battery capacity can be estimated as:

\begin{equation}
    C_{batt} = \frac{E_{night} \times D_{autonomy}}{DoD \times \eta_{batt}}
    \label{eq:battery_sizing}
\end{equation}

\noindent where:
\begin{itemize}
    \item $E_{night}$ = daily energy consumption during non-solar hours (Wh)
    \item $D_{autonomy}$ = desired days of autonomous operation (typically 1--2)
    \item $DoD$ = maximum depth of discharge (0.8 for LFP, 0.5 for lead-acid)
    \item $\eta_{batt}$ = battery round-trip efficiency (0.90--0.95 for Li-ion)
\end{itemize}

Research by Khezri et al.\ \cite{khezri2020} and Hesse et al.\ \cite{hesse2017} demonstrated that optimal sizing requires detailed household load profile analysis rather than generic rule-of-thumb approaches.


\section{Grid Connection and Net Metering}

\subsection{Grid Interconnection}

Grid interconnection enables bidirectional power flow between the PV system and the utility. Compliance with IEEE~1547 \cite{ieee1547} and IEC~61727 is required, specifying voltage/frequency operating ranges, harmonic limits, anti-islanding response, and grounding requirements.

\subsection{Net Metering}

Net metering credits PV system owners for surplus exported electricity using a bidirectional meter. In the Iraqi context, net metering policies are still in early development. The absence of mature net metering regulations currently limits purely on-grid system economics---further strengthening the case for hybrid systems with self-consumption optimization.


\section{System Architectures: On-Grid vs.\ Hybrid}

\subsection{On-Grid System Architecture}

An on-grid system consists of the PV array, grid-tied inverter, and balance-of-system components:

\begin{figure}[H]
    \centering
    \includegraphics[width=0.85\textwidth]{figures/fig2_5_safety_components.png}
    \caption{Safety components in a residential PV system including grounding, protection devices, DC disconnect, inverter, and emergency disconnect mechanisms.}
    \label{fig:safety}
\end{figure}

Key characteristics: no battery storage, system ceases during outages, lowest capital cost, highest conversion efficiency (no battery losses).


\subsection{Hybrid System Architecture}

Hybrid systems add battery storage and a hybrid inverter. Unlike standard grid-tied configurations, the hybrid inverter is programmed with a dynamic energy management algorithm that prioritizes power flows to maximize self-reliance and economic savings:

\begin{enumerate}
    \item \textbf{Supply immediate loads:} Available solar generation is first directed to satisfy immediate household electrical loads.
    \item \textbf{Charge the battery bank:} Any excess generation is subsequently routed to charge the batteries.
    \item \textbf{Export surplus power:} Only once the batteries reach a fully charged state is surplus power exported to the utility grid.
    \item \textbf{Discharge batteries:} During periods of insufficient solar irradiance or at night, the system defaults to discharging the batteries to meet household demand.
    \item \textbf{Import from utility:} The grid acts strictly as a tertiary source of last resort, supplying power only when both solar generation and battery reserves are depleted.
\end{enumerate}

In the event of a utility grid failure, an internal transfer switch rapidly isolates the system from the grid and transitions the inverter into island mode, allowing it to continue supplying critical household loads from the PV array and battery reserves---a decisive operational advantage in regions plagued by chronic infrastructure instability.

\begin{table}[H]
\centering
\caption{Comparative Summary: On-Grid vs.\ Hybrid System Architectures}
\label{tab:architecture_comparison}
\small
\begin{tabular}{@{}p{4cm} p{4.5cm} p{4.5cm}@{}}
\toprule
\textbf{Feature} & \textbf{On-Grid} & \textbf{Hybrid} \\
\midrule
Battery Storage & None & Yes (Li-ion or lead-acid) \\
Grid Dependency & High & Medium \\
Self-Consumption & 30--40\% & 60--80\% \\
Capital Cost & Lower & Higher (30--60\% premium) \\
Complexity & Low & Medium--High \\
Energy Independence & Low & Medium--High \\
During Outages & Shuts down & Operates (island mode) \\
Conversion Efficiency & Higher & Slightly lower (5--10\% battery losses) \\
Suitability for Erbil & Limited & Highly suitable \\
\bottomrule
\end{tabular}
\end{table}


\section{Chapter Summary}

This chapter examined the core components of residential PV systems: (1)~PV modules and the photovoltaic effect, (2)~inverter types and MPPT technology, (3)~battery energy storage technologies and sizing, (4)~grid interconnection and net metering, and (5)~on-grid versus hybrid architectures. These foundations support the design calculations and performance analysis in subsequent chapters.
 from your main document

\chapter{SYSTEM COMPONENTS AND WORKING PRINCIPLES}

\section{Overview}

The design and implementation of residential photovoltaic (PV) systems---whether on-grid or hybrid---require an in-depth understanding of the individual components, their operating principles, and how they interact. This chapter presents a detailed examination of the core components: solar PV modules, inverters, battery energy storage systems, and grid connection mechanisms.


\section{Solar Photovoltaic (PV) Modules}

\subsection{Operating Principle: The Photovoltaic Effect}

Solar PV modules convert incident solar radiation directly into electrical energy through the photovoltaic effect---a quantum mechanical process first observed by Becquerel in 1839 and first practically demonstrated by Bell Laboratories in 1954 \cite{chapin1954}.

The photovoltaic effect occurs within semiconductor materials---most commonly crystalline silicon (Si)---engineered to contain a p-n junction. When photons strike the semiconductor surface with energy greater than the material's bandgap energy (approximately 1.12~eV for silicon), they excite valence electrons into the conduction band, creating electron-hole pairs. The internal electric field at the p-n junction separates these charge carriers: electrons are swept toward the n-type (phosphorus-doped) layer, while holes migrate toward the p-type (boron-doped) layer. This charge separation establishes a potential difference across the cell, and when an external circuit is connected, current flows \cite{green2003,luque2011}.

\begin{figure}[H]
    \centering
    \includegraphics[width=0.75\textwidth]{figures/fig2_1_pv_conversion.png}
    \caption{The conversion of solar radiation to electricity in PV systems, showing the p-n junction, photon absorption, electron flow, and hole flow mechanisms.}
    \label{fig:pv_conversion}
\end{figure}

The performance of a PV cell is characterized by several key electrical parameters measured under Standard Test Conditions (STC: 1000~W/m$^{2}$ irradiance, 25$^{\circ}$C cell temperature, and an Air Mass of 1.5):

\begin{itemize}
    \item \textbf{Open-Circuit Voltage ($V_{oc}$):} This represents the maximum voltage available from a solar cell when no current is flowing, typically ranging from 0.5 to 0.7~V for standard silicon cells.
    \item \textbf{Short-Circuit Current ($I_{sc}$):} This is the maximum current output at zero voltage, which is directly proportional to the incident solar irradiance and the active cell area.
    \item \textbf{Maximum Power Point ($P_{max}$):} During normal operation, the cell operates at this specific voltage and current point where the electrical power output is maximized.
    \item \textbf{Fill Factor (FF):} Defined as the ratio of $P_{max}$ to the product of $V_{oc}$ and $I_{sc}$. High-quality cells typically exhibit a fill factor between 0.75 and 0.85.
    \item \textbf{Conversion Efficiency ($\eta$):} This dictates the proportion of incident solar power that is successfully converted into usable electrical power.
\end{itemize}

The fundamental relationships are:

\begin{equation}
    P_{max} = V_{mp} \times I_{mp} = FF \times V_{oc} \times I_{sc}
    \label{eq:pmax}
\end{equation}

\begin{equation}
    \eta = \frac{P_{max}}{G \times A} \times 100\%
    \label{eq:efficiency}
\end{equation}

\noindent where $G$ is solar irradiance (W/m$^{2}$) and $A$ is cell area (m$^{2}$).

\begin{figure}[H]
    \centering
    \includegraphics[width=0.7\textwidth]{figures/fig2_2_iv_curve.png}
    \caption{Typical I-V and P-V characteristic curves of a solar cell showing open-circuit voltage ($V_{oc}$), short-circuit current ($I_{sc}$), and the maximum power point ($P_{max}$).}
    \label{fig:iv_curve}
\end{figure}


\subsection{PV Module Construction and Types}

A PV module consists of multiple cells electrically connected in series and parallel to achieve desired voltage and current ratings. Cells are encapsulated between a front glass cover and rear backsheet with EVA encapsulant, providing mechanical protection and moisture resistance.

\begin{table}[H]
\centering
\caption{Comparison of Major PV Module Technologies}
\label{tab:pv_comparison}
\small
\begin{tabular}{@{}p{3.8cm} c c c@{}}
\toprule
\textbf{Parameter} & \textbf{Monocrystalline Si} & \textbf{Polycrystalline Si} & \textbf{Thin-Film} \\
\midrule
Cell Efficiency & 18--24\% & 16--20\% & 10--13\% \\
Module Efficiency & 17--22\% & 15--18\% & 9--12\% \\
Temp.\ Coefficient & --0.3 to --0.4 \%/$^{\circ}$C & --0.4 to --0.5 \%/$^{\circ}$C & --0.2 to --0.3 \%/$^{\circ}$C \\
Degradation Rate & 0.3--0.5 \%/yr & 0.5--0.7 \%/yr & 0.5--1.0 \%/yr \\
Lifespan & 25--30+ years & 25--30 years & 20--25 years \\
Space Efficiency & High & Medium & Low \\
Cost per Watt & Moderate--High & Low--Moderate & Low \\
\bottomrule
\end{tabular}
\end{table}

Monocrystalline silicon modules dominate the residential market due to superior efficiency and favorable temperature coefficients. The transition from PERC to TOPCon and HJT architectures has further improved efficiencies while reducing degradation \cite{longi2024,green2022}.


\subsection{Environmental Factors Affecting Performance}

The actual field output of a PV module is highly dependent on localized environmental parameters that cause deviations from STC ratings:

\begin{itemize}
    \item \textbf{Operating Temperature:} Module efficiency exhibits a negative correlation with increasing cell temperature, typically degrading at a rate of --0.3 to --0.5~\%/$^{\circ}$C. In Erbil's extreme summer climate, where rooftop module temperatures can reach 60--70$^{\circ}$C, temperature-induced power losses can reach 15--20\% relative to baseline STC values \cite{skoplaki2009}.
    \item \textbf{Soiling and Dust Accumulation:} The accumulation of airborne particulate matter poses a significant challenge. Studies in arid Middle Eastern climates indicate soiling can induce yield losses of 15--30\% if regular maintenance is neglected, making periodic cleaning essential during the Kurdistan Region's dry summers \cite{mani2010}.
    \item \textbf{Wind Speed:} While moderate wind velocities can marginally improve performance by providing convective cooling to the modules, extreme wind loads necessitate robust structural mounting to prevent mechanical failure.
    \item \textbf{Shading:} Even partial shading from nearby structures or vegetation can disproportionately reduce total array output. Because cells within a module are series-connected, a shaded cell acts as a reverse-biased resistor, restricting current flow and creating localized hot spots.
\end{itemize}

\begin{figure}[H]
    \centering
    \includegraphics[width=0.75\textwidth]{figures/fig2_3_rooftop_orientation.png}
    \caption{Typical residential rooftop solar PV installation showing panel arrangement, mounting structure, tilt angle, and exposure to sunlight with minimal shading conditions.}
    \label{fig:rooftop}
\end{figure}


\section{Inverters}

\subsection{Function and Operating Principle}

The primary function of a solar inverter is the conversion of direct current (DC) electricity generated by the PV array into alternating current (AC). However, modern inverters act as the central intelligence hub of the PV system, executing several critical ancillary functions:

\begin{itemize}
    \item \textbf{Maximum Power Point Tracking (MPPT):} The inverter continuously samples the voltage and current of the PV array and adjusts the electrical operating point to ensure maximum power extraction under fluctuating conditions, frequently achieving tracking efficiencies exceeding 99\% \cite{esram2007}.
    \item \textbf{Grid Synchronization:} For grid-tied operations, the inverter precisely matches the output voltage, phase, and frequency to the local utility specifications (50~Hz in Iraq) while maintaining Total Harmonic Distortion (THD) below 3--5\%.
    \item \textbf{Anti-Islanding Protection:} Crucially for safety, grid-interactive inverters must rapidly detect utility outages and automatically disconnect from the grid per IEEE~1547 standards \cite{ieee1547}, preventing power back-feed into compromised distribution lines.
    \item \textbf{System Monitoring:} Contemporary inverters integrate robust data logging and communication protocols (Wi-Fi, Ethernet, RS485), enabling real-time system monitoring, fault diagnostics, and remote management.
\end{itemize}


\subsection{Types of Inverters}

\begin{table}[H]
\centering
\caption{Comparison of Inverter Types for Residential PV Systems}
\label{tab:inverter_comparison}
\small
\begin{tabular}{@{}p{3.5cm} c c c@{}}
\toprule
\textbf{Parameter} & \textbf{String} & \textbf{Microinverter} & \textbf{Hybrid} \\
\midrule
Peak Efficiency & 96--98\% & 95--97\% & 93--97\% \\
MPPT Level & System (1--2) & Module-level & System + Battery \\
Shading Tolerance & Low & High & Low--Medium \\
Battery Integration & No & No & Yes (built-in) \\
Islanding Capability & No & No & Yes \\
Cost per Watt & Lowest & Highest & Moderate--High \\
Typical Lifespan & 10--15 yr & 15--25 yr & 10--15 yr \\
\bottomrule
\end{tabular}
\end{table}

For this study, a hybrid inverter is employed for the hybrid system, while a standard string inverter is utilized for the on-grid system.


\section{Battery Energy Storage Systems (BESS)}

\subsection{Role of Battery Storage}

Within a hybrid PV architecture, Battery Energy Storage Systems (BESS) are utilized to resolve the fundamental temporal mismatch between peak solar generation and residential consumption, serving three primary strategic functions:

\begin{itemize}
    \item \textbf{Energy Time-Shifting:} By storing excess diurnal generation for nocturnal use, households can increase their volumetric self-consumption of solar energy from a baseline of 30--40\% to upwards of 60--80\% \cite{koskela2019}.
    \item \textbf{Backup Power:} In the context of the Kurdistan Region's unreliable electricity infrastructure, batteries provide uninterruptible power supply during sudden utility grid outages.
    \item \textbf{Peak Shaving:} Intelligent hybrid inverters can utilize stored battery capacity to deliberately discharge during periods of high grid demand, minimizing reliance on expensive utility tariffs or neighborhood diesel generator subscriptions.
\end{itemize}

\begin{figure}[H]
    \centering
    \includegraphics[width=0.8\textwidth]{figures/fig2_4_wiring_diagram.png}
    \caption{Electrical wiring and connections in a residential PV system showing integration between solar panels, charge controller, battery, inverter, and load.}
    \label{fig:wiring}
\end{figure}


\subsection{Battery Technologies}

\begin{table}[H]
\centering
\caption{Comparison of Battery Technologies for Residential PV Systems}
\label{tab:battery_comparison}
\small
\begin{tabular}{@{}p{3.8cm} c c c@{}}
\toprule
\textbf{Parameter} & \textbf{Lead-Acid} & \textbf{Li-Ion (NMC)} & \textbf{LiFePO4 (LFP)} \\
\midrule
Energy Density & 30--50 Wh/kg & 150--220 Wh/kg & 90--160 Wh/kg \\
Cycle Life (80\% DoD) & 300--800 & 2,000--4,000 & 3,000--10,000 \\
Recommended DoD & 50\% & 80--90\% & 80--100\% \\
Round-Trip Efficiency & 75--85\% & 90--95\% & 92--98\% \\
Calendar Life & 3--5 yr & 8--12 yr & 10--15 yr \\
Safety & H$_2$ off-gassing & Thermal runaway risk & Excellent stability \\
Cost per kWh (2024) & \$100--200 & \$250--400 & \$200--350 \\
\bottomrule
\end{tabular}
\end{table}

LFP chemistry is increasingly preferred for residential applications due to excellent thermal stability, non-toxicity, and extended cycle life \cite{zakeri2015,hesse2017}.


\subsection{Battery Sizing}

The required battery capacity can be estimated as:

\begin{equation}
    C_{batt} = \frac{E_{night} \times D_{autonomy}}{DoD \times \eta_{batt}}
    \label{eq:battery_sizing}
\end{equation}

\noindent where:
\begin{itemize}
    \item $E_{night}$ = daily energy consumption during non-solar hours (Wh)
    \item $D_{autonomy}$ = desired days of autonomous operation (typically 1--2)
    \item $DoD$ = maximum depth of discharge (0.8 for LFP, 0.5 for lead-acid)
    \item $\eta_{batt}$ = battery round-trip efficiency (0.90--0.95 for Li-ion)
\end{itemize}

Research by Khezri et al.\ \cite{khezri2020} and Hesse et al.\ \cite{hesse2017} demonstrated that optimal sizing requires detailed household load profile analysis rather than generic rule-of-thumb approaches.


\section{Grid Connection and Net Metering}

\subsection{Grid Interconnection}

Grid interconnection enables bidirectional power flow between the PV system and the utility. Compliance with IEEE~1547 \cite{ieee1547} and IEC~61727 is required, specifying voltage/frequency operating ranges, harmonic limits, anti-islanding response, and grounding requirements.

\subsection{Net Metering}

Net metering credits PV system owners for surplus exported electricity using a bidirectional meter. In the Iraqi context, net metering policies are still in early development. The absence of mature net metering regulations currently limits purely on-grid system economics---further strengthening the case for hybrid systems with self-consumption optimization.


\section{System Architectures: On-Grid vs.\ Hybrid}

\subsection{On-Grid System Architecture}

An on-grid system consists of the PV array, grid-tied inverter, and balance-of-system components:

\begin{figure}[H]
    \centering
    \includegraphics[width=0.85\textwidth]{figures/fig2_5_safety_components.png}
    \caption{Safety components in a residential PV system including grounding, protection devices, DC disconnect, inverter, and emergency disconnect mechanisms.}
    \label{fig:safety}
\end{figure}

Key characteristics: no battery storage, system ceases during outages, lowest capital cost, highest conversion efficiency (no battery losses).


\subsection{Hybrid System Architecture}

Hybrid systems add battery storage and a hybrid inverter. Unlike standard grid-tied configurations, the hybrid inverter is programmed with a dynamic energy management algorithm that prioritizes power flows to maximize self-reliance and economic savings:

\begin{enumerate}
    \item \textbf{Supply immediate loads:} Available solar generation is first directed to satisfy immediate household electrical loads.
    \item \textbf{Charge the battery bank:} Any excess generation is subsequently routed to charge the batteries.
    \item \textbf{Export surplus power:} Only once the batteries reach a fully charged state is surplus power exported to the utility grid.
    \item \textbf{Discharge batteries:} During periods of insufficient solar irradiance or at night, the system defaults to discharging the batteries to meet household demand.
    \item \textbf{Import from utility:} The grid acts strictly as a tertiary source of last resort, supplying power only when both solar generation and battery reserves are depleted.
\end{enumerate}

In the event of a utility grid failure, an internal transfer switch rapidly isolates the system from the grid and transitions the inverter into island mode, allowing it to continue supplying critical household loads from the PV array and battery reserves---a decisive operational advantage in regions plagued by chronic infrastructure instability.

\begin{table}[H]
\centering
\caption{Comparative Summary: On-Grid vs.\ Hybrid System Architectures}
\label{tab:architecture_comparison}
\small
\begin{tabular}{@{}p{4cm} p{4.5cm} p{4.5cm}@{}}
\toprule
\textbf{Feature} & \textbf{On-Grid} & \textbf{Hybrid} \\
\midrule
Battery Storage & None & Yes (Li-ion or lead-acid) \\
Grid Dependency & High & Medium \\
Self-Consumption & 30--40\% & 60--80\% \\
Capital Cost & Lower & Higher (30--60\% premium) \\
Complexity & Low & Medium--High \\
Energy Independence & Low & Medium--High \\
During Outages & Shuts down & Operates (island mode) \\
Conversion Efficiency & Higher & Slightly lower (5--10\% battery losses) \\
Suitability for Erbil & Limited & Highly suitable \\
\bottomrule
\end{tabular}
\end{table}


\section{Chapter Summary}

This chapter examined the core components of residential PV systems: (1)~PV modules and the photovoltaic effect, (2)~inverter types and MPPT technology, (3)~battery energy storage technologies and sizing, (4)~grid interconnection and net metering, and (5)~on-grid versus hybrid architectures. These foundations support the design calculations and performance analysis in subsequent chapters.
 from your main document

\chapter{SYSTEM COMPONENTS AND WORKING PRINCIPLES}

\section{Overview}

The design and implementation of residential photovoltaic (PV) systems---whether on-grid or hybrid---require an in-depth understanding of the individual components, their operating principles, and how they interact. This chapter presents a detailed examination of the core components: solar PV modules, inverters, battery energy storage systems, and grid connection mechanisms.


\section{Solar Photovoltaic (PV) Modules}

\subsection{Operating Principle: The Photovoltaic Effect}

Solar PV modules convert incident solar radiation directly into electrical energy through the photovoltaic effect---a quantum mechanical process first observed by Becquerel in 1839 and first practically demonstrated by Bell Laboratories in 1954 \cite{chapin1954}.

The photovoltaic effect occurs within semiconductor materials---most commonly crystalline silicon (Si)---engineered to contain a p-n junction. When photons strike the semiconductor surface with energy greater than the material's bandgap energy (approximately 1.12~eV for silicon), they excite valence electrons into the conduction band, creating electron-hole pairs. The internal electric field at the p-n junction separates these charge carriers: electrons are swept toward the n-type (phosphorus-doped) layer, while holes migrate toward the p-type (boron-doped) layer. This charge separation establishes a potential difference across the cell, and when an external circuit is connected, current flows \cite{green2003,luque2011}.

\begin{figure}[H]
    \centering
    \includegraphics[width=0.75\textwidth]{figures/fig2_1_pv_conversion.png}
    \caption{The conversion of solar radiation to electricity in PV systems, showing the p-n junction, photon absorption, electron flow, and hole flow mechanisms.}
    \label{fig:pv_conversion}
\end{figure}

The performance of a PV cell is characterized by several key electrical parameters measured under Standard Test Conditions (STC: 1000~W/m$^{2}$ irradiance, 25$^{\circ}$C cell temperature, and an Air Mass of 1.5):

\begin{itemize}
    \item \textbf{Open-Circuit Voltage ($V_{oc}$):} This represents the maximum voltage available from a solar cell when no current is flowing, typically ranging from 0.5 to 0.7~V for standard silicon cells.
    \item \textbf{Short-Circuit Current ($I_{sc}$):} This is the maximum current output at zero voltage, which is directly proportional to the incident solar irradiance and the active cell area.
    \item \textbf{Maximum Power Point ($P_{max}$):} During normal operation, the cell operates at this specific voltage and current point where the electrical power output is maximized.
    \item \textbf{Fill Factor (FF):} Defined as the ratio of $P_{max}$ to the product of $V_{oc}$ and $I_{sc}$. High-quality cells typically exhibit a fill factor between 0.75 and 0.85.
    \item \textbf{Conversion Efficiency ($\eta$):} This dictates the proportion of incident solar power that is successfully converted into usable electrical power.
\end{itemize}

The fundamental relationships are:

\begin{equation}
    P_{max} = V_{mp} \times I_{mp} = FF \times V_{oc} \times I_{sc}
    \label{eq:pmax}
\end{equation}

\begin{equation}
    \eta = \frac{P_{max}}{G \times A} \times 100\%
    \label{eq:efficiency}
\end{equation}

\noindent where $G$ is solar irradiance (W/m$^{2}$) and $A$ is cell area (m$^{2}$).

\begin{figure}[H]
    \centering
    \includegraphics[width=0.7\textwidth]{figures/fig2_2_iv_curve.png}
    \caption{Typical I-V and P-V characteristic curves of a solar cell showing open-circuit voltage ($V_{oc}$), short-circuit current ($I_{sc}$), and the maximum power point ($P_{max}$).}
    \label{fig:iv_curve}
\end{figure}


\subsection{PV Module Construction and Types}

A PV module consists of multiple cells electrically connected in series and parallel to achieve desired voltage and current ratings. Cells are encapsulated between a front glass cover and rear backsheet with EVA encapsulant, providing mechanical protection and moisture resistance.

\begin{table}[H]
\centering
\caption{Comparison of Major PV Module Technologies}
\label{tab:pv_comparison}
\small
\begin{tabular}{@{}p{3.8cm} c c c@{}}
\toprule
\textbf{Parameter} & \textbf{Monocrystalline Si} & \textbf{Polycrystalline Si} & \textbf{Thin-Film} \\
\midrule
Cell Efficiency & 18--24\% & 16--20\% & 10--13\% \\
Module Efficiency & 17--22\% & 15--18\% & 9--12\% \\
Temp.\ Coefficient & --0.3 to --0.4 \%/$^{\circ}$C & --0.4 to --0.5 \%/$^{\circ}$C & --0.2 to --0.3 \%/$^{\circ}$C \\
Degradation Rate & 0.3--0.5 \%/yr & 0.5--0.7 \%/yr & 0.5--1.0 \%/yr \\
Lifespan & 25--30+ years & 25--30 years & 20--25 years \\
Space Efficiency & High & Medium & Low \\
Cost per Watt & Moderate--High & Low--Moderate & Low \\
\bottomrule
\end{tabular}
\end{table}

Monocrystalline silicon modules dominate the residential market due to superior efficiency and favorable temperature coefficients. The transition from PERC to TOPCon and HJT architectures has further improved efficiencies while reducing degradation \cite{longi2024,green2022}.


\subsection{Environmental Factors Affecting Performance}

The actual field output of a PV module is highly dependent on localized environmental parameters that cause deviations from STC ratings:

\begin{itemize}
    \item \textbf{Operating Temperature:} Module efficiency exhibits a negative correlation with increasing cell temperature, typically degrading at a rate of --0.3 to --0.5~\%/$^{\circ}$C. In Erbil's extreme summer climate, where rooftop module temperatures can reach 60--70$^{\circ}$C, temperature-induced power losses can reach 15--20\% relative to baseline STC values \cite{skoplaki2009}.
    \item \textbf{Soiling and Dust Accumulation:} The accumulation of airborne particulate matter poses a significant challenge. Studies in arid Middle Eastern climates indicate soiling can induce yield losses of 15--30\% if regular maintenance is neglected, making periodic cleaning essential during the Kurdistan Region's dry summers \cite{mani2010}.
    \item \textbf{Wind Speed:} While moderate wind velocities can marginally improve performance by providing convective cooling to the modules, extreme wind loads necessitate robust structural mounting to prevent mechanical failure.
    \item \textbf{Shading:} Even partial shading from nearby structures or vegetation can disproportionately reduce total array output. Because cells within a module are series-connected, a shaded cell acts as a reverse-biased resistor, restricting current flow and creating localized hot spots.
\end{itemize}

\begin{figure}[H]
    \centering
    \includegraphics[width=0.75\textwidth]{figures/fig2_3_rooftop_orientation.png}
    \caption{Typical residential rooftop solar PV installation showing panel arrangement, mounting structure, tilt angle, and exposure to sunlight with minimal shading conditions.}
    \label{fig:rooftop}
\end{figure}


\section{Inverters}

\subsection{Function and Operating Principle}

The primary function of a solar inverter is the conversion of direct current (DC) electricity generated by the PV array into alternating current (AC). However, modern inverters act as the central intelligence hub of the PV system, executing several critical ancillary functions:

\begin{itemize}
    \item \textbf{Maximum Power Point Tracking (MPPT):} The inverter continuously samples the voltage and current of the PV array and adjusts the electrical operating point to ensure maximum power extraction under fluctuating conditions, frequently achieving tracking efficiencies exceeding 99\% \cite{esram2007}.
    \item \textbf{Grid Synchronization:} For grid-tied operations, the inverter precisely matches the output voltage, phase, and frequency to the local utility specifications (50~Hz in Iraq) while maintaining Total Harmonic Distortion (THD) below 3--5\%.
    \item \textbf{Anti-Islanding Protection:} Crucially for safety, grid-interactive inverters must rapidly detect utility outages and automatically disconnect from the grid per IEEE~1547 standards \cite{ieee1547}, preventing power back-feed into compromised distribution lines.
    \item \textbf{System Monitoring:} Contemporary inverters integrate robust data logging and communication protocols (Wi-Fi, Ethernet, RS485), enabling real-time system monitoring, fault diagnostics, and remote management.
\end{itemize}


\subsection{Types of Inverters}

\begin{table}[H]
\centering
\caption{Comparison of Inverter Types for Residential PV Systems}
\label{tab:inverter_comparison}
\small
\begin{tabular}{@{}p{3.5cm} c c c@{}}
\toprule
\textbf{Parameter} & \textbf{String} & \textbf{Microinverter} & \textbf{Hybrid} \\
\midrule
Peak Efficiency & 96--98\% & 95--97\% & 93--97\% \\
MPPT Level & System (1--2) & Module-level & System + Battery \\
Shading Tolerance & Low & High & Low--Medium \\
Battery Integration & No & No & Yes (built-in) \\
Islanding Capability & No & No & Yes \\
Cost per Watt & Lowest & Highest & Moderate--High \\
Typical Lifespan & 10--15 yr & 15--25 yr & 10--15 yr \\
\bottomrule
\end{tabular}
\end{table}

For this study, a hybrid inverter is employed for the hybrid system, while a standard string inverter is utilized for the on-grid system.


\section{Battery Energy Storage Systems (BESS)}

\subsection{Role of Battery Storage}

Within a hybrid PV architecture, Battery Energy Storage Systems (BESS) are utilized to resolve the fundamental temporal mismatch between peak solar generation and residential consumption, serving three primary strategic functions:

\begin{itemize}
    \item \textbf{Energy Time-Shifting:} By storing excess diurnal generation for nocturnal use, households can increase their volumetric self-consumption of solar energy from a baseline of 30--40\% to upwards of 60--80\% \cite{koskela2019}.
    \item \textbf{Backup Power:} In the context of the Kurdistan Region's unreliable electricity infrastructure, batteries provide uninterruptible power supply during sudden utility grid outages.
    \item \textbf{Peak Shaving:} Intelligent hybrid inverters can utilize stored battery capacity to deliberately discharge during periods of high grid demand, minimizing reliance on expensive utility tariffs or neighborhood diesel generator subscriptions.
\end{itemize}

\begin{figure}[H]
    \centering
    \includegraphics[width=0.8\textwidth]{figures/fig2_4_wiring_diagram.png}
    \caption{Electrical wiring and connections in a residential PV system showing integration between solar panels, charge controller, battery, inverter, and load.}
    \label{fig:wiring}
\end{figure}


\subsection{Battery Technologies}

\begin{table}[H]
\centering
\caption{Comparison of Battery Technologies for Residential PV Systems}
\label{tab:battery_comparison}
\small
\begin{tabular}{@{}p{3.8cm} c c c@{}}
\toprule
\textbf{Parameter} & \textbf{Lead-Acid} & \textbf{Li-Ion (NMC)} & \textbf{LiFePO4 (LFP)} \\
\midrule
Energy Density & 30--50 Wh/kg & 150--220 Wh/kg & 90--160 Wh/kg \\
Cycle Life (80\% DoD) & 300--800 & 2,000--4,000 & 3,000--10,000 \\
Recommended DoD & 50\% & 80--90\% & 80--100\% \\
Round-Trip Efficiency & 75--85\% & 90--95\% & 92--98\% \\
Calendar Life & 3--5 yr & 8--12 yr & 10--15 yr \\
Safety & H$_2$ off-gassing & Thermal runaway risk & Excellent stability \\
Cost per kWh (2024) & \$100--200 & \$250--400 & \$200--350 \\
\bottomrule
\end{tabular}
\end{table}

LFP chemistry is increasingly preferred for residential applications due to excellent thermal stability, non-toxicity, and extended cycle life \cite{zakeri2015,hesse2017}.


\subsection{Battery Sizing}

The required battery capacity can be estimated as:

\begin{equation}
    C_{batt} = \frac{E_{night} \times D_{autonomy}}{DoD \times \eta_{batt}}
    \label{eq:battery_sizing}
\end{equation}

\noindent where:
\begin{itemize}
    \item $E_{night}$ = daily energy consumption during non-solar hours (Wh)
    \item $D_{autonomy}$ = desired days of autonomous operation (typically 1--2)
    \item $DoD$ = maximum depth of discharge (0.8 for LFP, 0.5 for lead-acid)
    \item $\eta_{batt}$ = battery round-trip efficiency (0.90--0.95 for Li-ion)
\end{itemize}

Research by Khezri et al.\ \cite{khezri2020} and Hesse et al.\ \cite{hesse2017} demonstrated that optimal sizing requires detailed household load profile analysis rather than generic rule-of-thumb approaches.


\section{Grid Connection and Net Metering}

\subsection{Grid Interconnection}

Grid interconnection enables bidirectional power flow between the PV system and the utility. Compliance with IEEE~1547 \cite{ieee1547} and IEC~61727 is required, specifying voltage/frequency operating ranges, harmonic limits, anti-islanding response, and grounding requirements.

\subsection{Net Metering}

Net metering credits PV system owners for surplus exported electricity using a bidirectional meter. In the Iraqi context, net metering policies are still in early development. The absence of mature net metering regulations currently limits purely on-grid system economics---further strengthening the case for hybrid systems with self-consumption optimization.


\section{System Architectures: On-Grid vs.\ Hybrid}

\subsection{On-Grid System Architecture}

An on-grid system consists of the PV array, grid-tied inverter, and balance-of-system components:

\begin{figure}[H]
    \centering
    \includegraphics[width=0.85\textwidth]{figures/fig2_5_safety_components.png}
    \caption{Safety components in a residential PV system including grounding, protection devices, DC disconnect, inverter, and emergency disconnect mechanisms.}
    \label{fig:safety}
\end{figure}

Key characteristics: no battery storage, system ceases during outages, lowest capital cost, highest conversion efficiency (no battery losses).


\subsection{Hybrid System Architecture}

Hybrid systems add battery storage and a hybrid inverter. Unlike standard grid-tied configurations, the hybrid inverter is programmed with a dynamic energy management algorithm that prioritizes power flows to maximize self-reliance and economic savings:

\begin{enumerate}
    \item \textbf{Supply immediate loads:} Available solar generation is first directed to satisfy immediate household electrical loads.
    \item \textbf{Charge the battery bank:} Any excess generation is subsequently routed to charge the batteries.
    \item \textbf{Export surplus power:} Only once the batteries reach a fully charged state is surplus power exported to the utility grid.
    \item \textbf{Discharge batteries:} During periods of insufficient solar irradiance or at night, the system defaults to discharging the batteries to meet household demand.
    \item \textbf{Import from utility:} The grid acts strictly as a tertiary source of last resort, supplying power only when both solar generation and battery reserves are depleted.
\end{enumerate}

In the event of a utility grid failure, an internal transfer switch rapidly isolates the system from the grid and transitions the inverter into island mode, allowing it to continue supplying critical household loads from the PV array and battery reserves---a decisive operational advantage in regions plagued by chronic infrastructure instability.

\begin{table}[H]
\centering
\caption{Comparative Summary: On-Grid vs.\ Hybrid System Architectures}
\label{tab:architecture_comparison}
\small
\begin{tabular}{@{}p{4cm} p{4.5cm} p{4.5cm}@{}}
\toprule
\textbf{Feature} & \textbf{On-Grid} & \textbf{Hybrid} \\
\midrule
Battery Storage & None & Yes (Li-ion or lead-acid) \\
Grid Dependency & High & Medium \\
Self-Consumption & 30--40\% & 60--80\% \\
Capital Cost & Lower & Higher (30--60\% premium) \\
Complexity & Low & Medium--High \\
Energy Independence & Low & Medium--High \\
During Outages & Shuts down & Operates (island mode) \\
Conversion Efficiency & Higher & Slightly lower (5--10\% battery losses) \\
Suitability for Erbil & Limited & Highly suitable \\
\bottomrule
\end{tabular}
\end{table}


\section{Chapter Summary}

This chapter examined the core components of residential PV systems: (1)~PV modules and the photovoltaic effect, (2)~inverter types and MPPT technology, (3)~battery energy storage technologies and sizing, (4)~grid interconnection and net metering, and (5)~on-grid versus hybrid architectures. These foundations support the design calculations and performance analysis in subsequent chapters.
 from your main document

\chapter{SYSTEM COMPONENTS AND WORKING PRINCIPLES}

\section{Overview}

The design and implementation of residential photovoltaic (PV) systems---whether on-grid or hybrid---require an in-depth understanding of the individual components, their operating principles, and how they interact. This chapter presents a detailed examination of the core components: solar PV modules, inverters, battery energy storage systems, and grid connection mechanisms.


\section{Solar Photovoltaic (PV) Modules}

\subsection{Operating Principle: The Photovoltaic Effect}

Solar PV modules convert incident solar radiation directly into electrical energy through the photovoltaic effect---a quantum mechanical process first observed by Becquerel in 1839 and first practically demonstrated by Bell Laboratories in 1954 \cite{chapin1954}.

The photovoltaic effect occurs within semiconductor materials---most commonly crystalline silicon (Si)---engineered to contain a p-n junction. When photons strike the semiconductor surface with energy greater than the material's bandgap energy (approximately 1.12~eV for silicon), they excite valence electrons into the conduction band, creating electron-hole pairs. The internal electric field at the p-n junction separates these charge carriers: electrons are swept toward the n-type (phosphorus-doped) layer, while holes migrate toward the p-type (boron-doped) layer. This charge separation establishes a potential difference across the cell, and when an external circuit is connected, current flows \cite{green2003,luque2011}.

\begin{figure}[H]
    \centering
    \includegraphics[width=0.75\textwidth]{figures/fig2_1_pv_conversion.png}
    \caption{The conversion of solar radiation to electricity in PV systems, showing the p-n junction, photon absorption, electron flow, and hole flow mechanisms.}
    \label{fig:pv_conversion}
\end{figure}

The performance of a PV cell is characterized by several key electrical parameters measured under Standard Test Conditions (STC: 1000~W/m$^{2}$ irradiance, 25$^{\circ}$C cell temperature, and an Air Mass of 1.5):

\begin{itemize}
    \item \textbf{Open-Circuit Voltage ($V_{oc}$):} This represents the maximum voltage available from a solar cell when no current is flowing, typically ranging from 0.5 to 0.7~V for standard silicon cells.
    \item \textbf{Short-Circuit Current ($I_{sc}$):} This is the maximum current output at zero voltage, which is directly proportional to the incident solar irradiance and the active cell area.
    \item \textbf{Maximum Power Point ($P_{max}$):} During normal operation, the cell operates at this specific voltage and current point where the electrical power output is maximized.
    \item \textbf{Fill Factor (FF):} Defined as the ratio of $P_{max}$ to the product of $V_{oc}$ and $I_{sc}$. High-quality cells typically exhibit a fill factor between 0.75 and 0.85.
    \item \textbf{Conversion Efficiency ($\eta$):} This dictates the proportion of incident solar power that is successfully converted into usable electrical power.
\end{itemize}

The fundamental relationships are:

\begin{equation}
    P_{max} = V_{mp} \times I_{mp} = FF \times V_{oc} \times I_{sc}
    \label{eq:pmax}
\end{equation}

\begin{equation}
    \eta = \frac{P_{max}}{G \times A} \times 100\%
    \label{eq:efficiency}
\end{equation}

\noindent where $G$ is solar irradiance (W/m$^{2}$) and $A$ is cell area (m$^{2}$).

\begin{figure}[H]
    \centering
    \includegraphics[width=0.7\textwidth]{figures/fig2_2_iv_curve.png}
    \caption{Typical I-V and P-V characteristic curves of a solar cell showing open-circuit voltage ($V_{oc}$), short-circuit current ($I_{sc}$), and the maximum power point ($P_{max}$).}
    \label{fig:iv_curve}
\end{figure}


\subsection{PV Module Construction and Types}

A PV module consists of multiple cells electrically connected in series and parallel to achieve desired voltage and current ratings. Cells are encapsulated between a front glass cover and rear backsheet with EVA encapsulant, providing mechanical protection and moisture resistance.

\begin{table}[H]
\centering
\caption{Comparison of Major PV Module Technologies}
\label{tab:pv_comparison}
\small
\begin{tabular}{@{}p{3.8cm} c c c@{}}
\toprule
\textbf{Parameter} & \textbf{Monocrystalline Si} & \textbf{Polycrystalline Si} & \textbf{Thin-Film} \\
\midrule
Cell Efficiency & 18--24\% & 16--20\% & 10--13\% \\
Module Efficiency & 17--22\% & 15--18\% & 9--12\% \\
Temp.\ Coefficient & --0.3 to --0.4 \%/$^{\circ}$C & --0.4 to --0.5 \%/$^{\circ}$C & --0.2 to --0.3 \%/$^{\circ}$C \\
Degradation Rate & 0.3--0.5 \%/yr & 0.5--0.7 \%/yr & 0.5--1.0 \%/yr \\
Lifespan & 25--30+ years & 25--30 years & 20--25 years \\
Space Efficiency & High & Medium & Low \\
Cost per Watt & Moderate--High & Low--Moderate & Low \\
\bottomrule
\end{tabular}
\end{table}

Monocrystalline silicon modules dominate the residential market due to superior efficiency and favorable temperature coefficients. The transition from PERC to TOPCon and HJT architectures has further improved efficiencies while reducing degradation \cite{longi2024,green2022}.


\subsection{Environmental Factors Affecting Performance}

The actual field output of a PV module is highly dependent on localized environmental parameters that cause deviations from STC ratings:

\begin{itemize}
    \item \textbf{Operating Temperature:} Module efficiency exhibits a negative correlation with increasing cell temperature, typically degrading at a rate of --0.3 to --0.5~\%/$^{\circ}$C. In Erbil's extreme summer climate, where rooftop module temperatures can reach 60--70$^{\circ}$C, temperature-induced power losses can reach 15--20\% relative to baseline STC values \cite{skoplaki2009}.
    \item \textbf{Soiling and Dust Accumulation:} The accumulation of airborne particulate matter poses a significant challenge. Studies in arid Middle Eastern climates indicate soiling can induce yield losses of 15--30\% if regular maintenance is neglected, making periodic cleaning essential during the Kurdistan Region's dry summers \cite{mani2010}.
    \item \textbf{Wind Speed:} While moderate wind velocities can marginally improve performance by providing convective cooling to the modules, extreme wind loads necessitate robust structural mounting to prevent mechanical failure.
    \item \textbf{Shading:} Even partial shading from nearby structures or vegetation can disproportionately reduce total array output. Because cells within a module are series-connected, a shaded cell acts as a reverse-biased resistor, restricting current flow and creating localized hot spots.
\end{itemize}

\begin{figure}[H]
    \centering
    \includegraphics[width=0.75\textwidth]{figures/fig2_3_rooftop_orientation.png}
    \caption{Typical residential rooftop solar PV installation showing panel arrangement, mounting structure, tilt angle, and exposure to sunlight with minimal shading conditions.}
    \label{fig:rooftop}
\end{figure}


\section{Inverters}

\subsection{Function and Operating Principle}

The primary function of a solar inverter is the conversion of direct current (DC) electricity generated by the PV array into alternating current (AC). However, modern inverters act as the central intelligence hub of the PV system, executing several critical ancillary functions:

\begin{itemize}
    \item \textbf{Maximum Power Point Tracking (MPPT):} The inverter continuously samples the voltage and current of the PV array and adjusts the electrical operating point to ensure maximum power extraction under fluctuating conditions, frequently achieving tracking efficiencies exceeding 99\% \cite{esram2007}.
    \item \textbf{Grid Synchronization:} For grid-tied operations, the inverter precisely matches the output voltage, phase, and frequency to the local utility specifications (50~Hz in Iraq) while maintaining Total Harmonic Distortion (THD) below 3--5\%.
    \item \textbf{Anti-Islanding Protection:} Crucially for safety, grid-interactive inverters must rapidly detect utility outages and automatically disconnect from the grid per IEEE~1547 standards \cite{ieee1547}, preventing power back-feed into compromised distribution lines.
    \item \textbf{System Monitoring:} Contemporary inverters integrate robust data logging and communication protocols (Wi-Fi, Ethernet, RS485), enabling real-time system monitoring, fault diagnostics, and remote management.
\end{itemize}


\subsection{Types of Inverters}

\begin{table}[H]
\centering
\caption{Comparison of Inverter Types for Residential PV Systems}
\label{tab:inverter_comparison}
\small
\begin{tabular}{@{}p{3.5cm} c c c@{}}
\toprule
\textbf{Parameter} & \textbf{String} & \textbf{Microinverter} & \textbf{Hybrid} \\
\midrule
Peak Efficiency & 96--98\% & 95--97\% & 93--97\% \\
MPPT Level & System (1--2) & Module-level & System + Battery \\
Shading Tolerance & Low & High & Low--Medium \\
Battery Integration & No & No & Yes (built-in) \\
Islanding Capability & No & No & Yes \\
Cost per Watt & Lowest & Highest & Moderate--High \\
Typical Lifespan & 10--15 yr & 15--25 yr & 10--15 yr \\
\bottomrule
\end{tabular}
\end{table}

For this study, a hybrid inverter is employed for the hybrid system, while a standard string inverter is utilized for the on-grid system.


\section{Battery Energy Storage Systems (BESS)}

\subsection{Role of Battery Storage}

Within a hybrid PV architecture, Battery Energy Storage Systems (BESS) are utilized to resolve the fundamental temporal mismatch between peak solar generation and residential consumption, serving three primary strategic functions:

\begin{itemize}
    \item \textbf{Energy Time-Shifting:} By storing excess diurnal generation for nocturnal use, households can increase their volumetric self-consumption of solar energy from a baseline of 30--40\% to upwards of 60--80\% \cite{koskela2019}.
    \item \textbf{Backup Power:} In the context of the Kurdistan Region's unreliable electricity infrastructure, batteries provide uninterruptible power supply during sudden utility grid outages.
    \item \textbf{Peak Shaving:} Intelligent hybrid inverters can utilize stored battery capacity to deliberately discharge during periods of high grid demand, minimizing reliance on expensive utility tariffs or neighborhood diesel generator subscriptions.
\end{itemize}

\begin{figure}[H]
    \centering
    \includegraphics[width=0.8\textwidth]{figures/fig2_4_wiring_diagram.png}
    \caption{Electrical wiring and connections in a residential PV system showing integration between solar panels, charge controller, battery, inverter, and load.}
    \label{fig:wiring}
\end{figure}


\subsection{Battery Technologies}

\begin{table}[H]
\centering
\caption{Comparison of Battery Technologies for Residential PV Systems}
\label{tab:battery_comparison}
\small
\begin{tabular}{@{}p{3.8cm} c c c@{}}
\toprule
\textbf{Parameter} & \textbf{Lead-Acid} & \textbf{Li-Ion (NMC)} & \textbf{LiFePO4 (LFP)} \\
\midrule
Energy Density & 30--50 Wh/kg & 150--220 Wh/kg & 90--160 Wh/kg \\
Cycle Life (80\% DoD) & 300--800 & 2,000--4,000 & 3,000--10,000 \\
Recommended DoD & 50\% & 80--90\% & 80--100\% \\
Round-Trip Efficiency & 75--85\% & 90--95\% & 92--98\% \\
Calendar Life & 3--5 yr & 8--12 yr & 10--15 yr \\
Safety & H$_2$ off-gassing & Thermal runaway risk & Excellent stability \\
Cost per kWh (2024) & \$100--200 & \$250--400 & \$200--350 \\
\bottomrule
\end{tabular}
\end{table}

LFP chemistry is increasingly preferred for residential applications due to excellent thermal stability, non-toxicity, and extended cycle life \cite{zakeri2015,hesse2017}.


\subsection{Battery Sizing}

The required battery capacity can be estimated as:

\begin{equation}
    C_{batt} = \frac{E_{night} \times D_{autonomy}}{DoD \times \eta_{batt}}
    \label{eq:battery_sizing}
\end{equation}

\noindent where:
\begin{itemize}
    \item $E_{night}$ = daily energy consumption during non-solar hours (Wh)
    \item $D_{autonomy}$ = desired days of autonomous operation (typically 1--2)
    \item $DoD$ = maximum depth of discharge (0.8 for LFP, 0.5 for lead-acid)
    \item $\eta_{batt}$ = battery round-trip efficiency (0.90--0.95 for Li-ion)
\end{itemize}

Research by Khezri et al.\ \cite{khezri2020} and Hesse et al.\ \cite{hesse2017} demonstrated that optimal sizing requires detailed household load profile analysis rather than generic rule-of-thumb approaches.


\section{Grid Connection and Net Metering}

\subsection{Grid Interconnection}

Grid interconnection enables bidirectional power flow between the PV system and the utility. Compliance with IEEE~1547 \cite{ieee1547} and IEC~61727 is required, specifying voltage/frequency operating ranges, harmonic limits, anti-islanding response, and grounding requirements.

\subsection{Net Metering}

Net metering credits PV system owners for surplus exported electricity using a bidirectional meter. In the Iraqi context, net metering policies are still in early development. The absence of mature net metering regulations currently limits purely on-grid system economics---further strengthening the case for hybrid systems with self-consumption optimization.


\section{System Architectures: On-Grid vs.\ Hybrid}

\subsection{On-Grid System Architecture}

An on-grid system consists of the PV array, grid-tied inverter, and balance-of-system components:

\begin{figure}[H]
    \centering
    \includegraphics[width=0.85\textwidth]{figures/fig2_5_safety_components.png}
    \caption{Safety components in a residential PV system including grounding, protection devices, DC disconnect, inverter, and emergency disconnect mechanisms.}
    \label{fig:safety}
\end{figure}

Key characteristics: no battery storage, system ceases during outages, lowest capital cost, highest conversion efficiency (no battery losses).


\subsection{Hybrid System Architecture}

Hybrid systems add battery storage and a hybrid inverter. Unlike standard grid-tied configurations, the hybrid inverter is programmed with a dynamic energy management algorithm that prioritizes power flows to maximize self-reliance and economic savings:

\begin{enumerate}
    \item \textbf{Supply immediate loads:} Available solar generation is first directed to satisfy immediate household electrical loads.
    \item \textbf{Charge the battery bank:} Any excess generation is subsequently routed to charge the batteries.
    \item \textbf{Export surplus power:} Only once the batteries reach a fully charged state is surplus power exported to the utility grid.
    \item \textbf{Discharge batteries:} During periods of insufficient solar irradiance or at night, the system defaults to discharging the batteries to meet household demand.
    \item \textbf{Import from utility:} The grid acts strictly as a tertiary source of last resort, supplying power only when both solar generation and battery reserves are depleted.
\end{enumerate}

In the event of a utility grid failure, an internal transfer switch rapidly isolates the system from the grid and transitions the inverter into island mode, allowing it to continue supplying critical household loads from the PV array and battery reserves---a decisive operational advantage in regions plagued by chronic infrastructure instability.

\begin{table}[H]
\centering
\caption{Comparative Summary: On-Grid vs.\ Hybrid System Architectures}
\label{tab:architecture_comparison}
\small
\begin{tabular}{@{}p{4cm} p{4.5cm} p{4.5cm}@{}}
\toprule
\textbf{Feature} & \textbf{On-Grid} & \textbf{Hybrid} \\
\midrule
Battery Storage & None & Yes (Li-ion or lead-acid) \\
Grid Dependency & High & Medium \\
Self-Consumption & 30--40\% & 60--80\% \\
Capital Cost & Lower & Higher (30--60\% premium) \\
Complexity & Low & Medium--High \\
Energy Independence & Low & Medium--High \\
During Outages & Shuts down & Operates (island mode) \\
Conversion Efficiency & Higher & Slightly lower (5--10\% battery losses) \\
Suitability for Erbil & Limited & Highly suitable \\
\bottomrule
\end{tabular}
\end{table}


\section{Chapter Summary}

This chapter examined the core components of residential PV systems: (1)~PV modules and the photovoltaic effect, (2)~inverter types and MPPT technology, (3)~battery energy storage technologies and sizing, (4)~grid interconnection and net metering, and (5)~on-grid versus hybrid architectures. These foundations support the design calculations and performance analysis in subsequent chapters.
